\scanpage{171}
\margin{170}
\uline{Réédition des E. V. T., chap. II} \mspage{(7)}
\margin{Tribu n\textsuperscript{o} 63}

l. 7, remplacer $V_\alpha$ par $g_\alpha(V_\alpha)$ ).

P. 55, vider les l. 1-8 du bas.

P. 56, vider les l. 1-2, et 1-12 du bas.

P. 57, l. 1 : II. Limites inductives d'espaces localement convexes. Se donner les $E_\alpha\subset E$ filtrants croissants avec $\bigcup E_\alpha=E$, avec les $E_\beta\longrightarrow E_\alpha$ continues.

Vider les l. 9-14.

l. 13 du bas : Exemple.

P. 58, l. 1, vider 1).

Vider les l. 1-13 du bas.

P. 59, après la déf. 2, donner les voisinages de 0.

(l. 1-8 du bas ; l. 8 du bas, attention au style).

Puis, dire que, pour tout $F$ localement convexe, on a $\mathcal{L}(E,F)=\prod\mathcal{L}(E_i,F)$, et que ceci caractérise la somme directe.

Corollaire : Une limite inductive est quotient d'une somme directe.

Proposition 6. L'injection canonique de $E$ dans $\prod E_i$ est continue. Si $I$ est fini, c'est un isomorphisme d'E. V. T.

P. 60, vider les l. 1-13 du bas.

P. 61, vider les l. 1-5.

Corollaire 1 (ex. 2) : Si les $G_i$ sont fermés, $\sum G_i$ est fermé.

Démonstration : $p_\alpha:E\longrightarrow E_\alpha$ est continu, donc $p_\alpha^{-1}(G_\alpha)$ est fermé, donc $\bigcap p_\alpha^{-1}(G_\alpha)=\sum G_\alpha$ est fermé.

Vider les p. 62, 63, 64 (l. 1-4). On annoncera quelque part qu'une somme directe de complets est complète, mais on le démontrera plus tard à l'aide du th. de Grothendieck.

P. 64, le cor. 1 disparaît, puisqu'il a été glorifié tout de suite après la déf. des sommes directes. Le cor. 2 sera bloqué avec la prop. 7 (dont on conservera les notations), et sera à peu près évident sur la définition catégorique des sommes directes.

Vider la Remarque 2.

P. 65, l. 12, il faut pouvoir dire que $E$ est la limite inductive des $E_n$.

Annoncer : Dans tout le n\textsuperscript{o}, $E$ sera la limite inductive des $E_n$.
